% =============================================================================
% Definicao de requisitos tecnicos e especificacoes-meta (projeto Riachuelo).
% Os codigos das vozes do cliente (N, D, DE) sao os mesmos da Secao 7
% (Tabela tab:matriz_ndd). As matrizes do QFD foram transcritas das imagens
% originais (imagens/definicao dos requisitos/fig01 e fig02) sem alterar
% nenhum valor; so os codigos das vozes e os nomes dos requisitos mudaram.
% =============================================================================

\section{DEFINIÇÃO DE REQUISITOS TÉCNICOS E ESPECIFICAÇÕES-META}
\label{sec:definicao_requisitos}

Para transformar as necessidades identificadas na Seção~\ref{sec:ndd} em características que orientem o desenvolvimento e a avaliação do produto, a equipe utilizou o método \textit{Quality Function Deployment} (QFD), que relaciona a voz do cliente a requisitos técnicos e evita que as necessidades sejam traduzidas diretamente em soluções sem análise técnica. Como o projeto será desenvolvido em escala de protótipo, as especificações precisaram ser compatíveis com os recursos financeiros, tecnológicos e materiais da equipe. O objetivo foi, portanto, definir critérios que permitam verificar se o conceito resolve as principais dores identificadas na pesquisa, sem fixar requisitos de uma solução industrial definitiva.

\subsection{Requisitos dos clientes}
\label{subsec:requisitos_dos_clientes}

A pesquisa com consumidores mostrou perfis com comportamentos distintos em relação à autonomia, ao atendimento humano, ao uso de tecnologia e à busca por produtos, sem que essas características se excluíssem. Um mesmo consumidor pode valorizar autonomia e rapidez em uma compra simples e preferir atendimento em uma compra com maior incerteza, como a de uma peça de maior valor ou a que exige uma segunda opinião. Por isso, a equipe não partiu da premissa de que substituir o atendimento humano por tecnologia seria desejável em todas as situações.

As entrevistas também registraram dificuldades de disponibilidade e localização de produtos, casos em que o consumidor precisou de um vendedor para verificar o estoque e episódios de perda de confiança em informações digitais de disponibilidade que não correspondiam à loja. No pagamento, alguns entrevistados relataram desistências de compra por causa de fila, sobretudo em períodos de maior movimento, enquanto outros demonstraram interesse em realizar etapas da compra sem interação constante com vendedores, desde que houvesse suporte disponível caso algo desse errado. Esses achados apontaram para uma solução que combine autonomia e rapidez com segurança e possibilidade de atendimento.

A partir desse conjunto, a equipe selecionou sete das vinte vozes da Seção~\ref{sec:ndd} para o QFD, de modo a representar as três categorias da análise (necessidades, desejos e demandas) e a priorizar as vozes relacionadas às funções que a integração entre a etiqueta antifurto e o aplicativo pode contemplar. Os códigos utilizados a seguir são os mesmos da Tabela~\ref{tab:matriz_ndd}.

Entre as necessidades, foram selecionadas quatro. A primeira, \textbf{Encontrar o produto procurado com facilidade (N03)}, refere-se à etapa de busca, em que a dificuldade pode impedir a compra mesmo quando há interesse pelo produto. A evidência concentrou-se no Cluster 3, no qual a dificuldade foi relatada de forma intensa, e o aplicativo poderá apoiar a identificação do produto e o acesso a informações sobre a peça. A segunda, \textbf{Confiar na precisão do estoque informado (N05)}, é condição para que uma solução digital agregue valor: como as entrevistas registraram perda de confiança após divergências entre o site e a loja, disponibilizar informações sobre produtos não bastará se o usuário não confiar em sua precisão.

A terceira necessidade, \textbf{Contar com atendimento humano quando necessário (N08)}, reflete o achado de que autonomia tecnológica não implica rejeição ao atendimento. A solução deverá, assim, complementar o atendimento existente, para que o uso da etiqueta e do aplicativo não gere insegurança ou abandono da compra entre consumidores menos confortáveis com o autoatendimento. A quarta, \textbf{Concluir a compra sem constrangimentos (N06)}, liga-se diretamente à mudança que o produto introduz no procedimento de saída da loja. Embora o survey tenha mostrado pouca variação nesse aspecto, as entrevistas registraram o receio de ser visto como alguém que não pagou e o desconforto com verificações adicionais. Além de garantir que a transação ocorra corretamente, a solução deverá, portanto, transmitir segurança e legitimidade ao consumidor durante a saída.

Entre os desejos, a equipe selecionou \textbf{Receber confirmação imediata do pagamento (D03)}, pela relação direta com a experiência de pagamento e com a confiança do consumidor ao deixar a loja. As entrevistas indicaram a expectativa de um aviso inequívoco de que a operação foi concluída e de que o cliente está autorizado a sair com o produto. O item foi classificado como desejo porque aumenta o conforto e a segurança percebida sem ter aparecido como barreira universal à compra.

Por fim, duas demandas foram selecionadas. \textbf{Reduzir o tempo de espera no caixa (DE01)} está associada a um comportamento observável de abandono: o survey mostrou maior incidência de desistência por fila no Cluster 2, e as entrevistas trouxeram episódios concretos em promoções e no horário de almoço. \textbf{Ter pagamento rápido e atendimento adequado (DE05)}, sintetiza a variação do nível de autonomia desejado conforme o contexto: alguns entrevistados separaram a reposição rápida, em que valorizam agilidade, das compras mais complexas, em que preferem o apoio de um vendedor. A proposta deverá, então, permitir uma jornada rápida e autônoma quando o consumidor quiser, mantendo o atendimento disponível nos casos em que a assistência for valorizada.

As necessidades selecionadas concentram-se na obtenção de informações confiáveis, na busca pelo produto, na disponibilidade de suporte humano e na saída sem constrangimento. O desejo refere-se à qualidade da experiência de pagamento, e as demandas, a comportamentos concretos de compra, como o abandono diante de filas e a busca por agilidade em compras simples.

\subsection{Requisitos do produto}
\label{subsec:requisitos_do_produto}

Com base nas sete vozes selecionadas, a equipe definiu sete requisitos técnicos para o sistema formado pela etiqueta antifurto e pelo aplicativo. Cada requisito foi formulado primeiro em termos funcionais, isto é, pelo que o sistema deverá fazer, e só depois associado a um critério de desempenho. Um valor numérico foi fixado apenas quando havia justificativa para ele; nos demais casos, o critério será definido nos testes do protótipo, a partir de medições de referência. Os requisitos e seus critérios estão reunidos na Tabela~\ref{tab:requisitos}.

\begin{table}[htbp]
    \centering
    \caption{Requisitos técnicos e critérios de desempenho}
    \label{tab:requisitos}
    \footnotesize
    \renewcommand{\arraystretch}{1.3}
    \begin{tabularx}{\textwidth}{@{}
        >{\RaggedRight\arraybackslash}p{1.0cm}
        >{\RaggedRight\arraybackslash}p{2.8cm}
        >{\RaggedRight\arraybackslash}X
        >{\RaggedRight\arraybackslash}X@{}}
        \toprule
        \textbf{Código} & \textbf{Requisito} & \textbf{Formulação funcional} & \textbf{Critério de desempenho} \\
        \midrule
        R01 & Apoio à localização da peça na loja
            & O sistema deverá ajudar o cliente a identificar em que área da loja se encontra a peça procurada.
            & A especificação-meta é que pelo menos \textbf{90\% das solicitações apresentem a localização correta do produto} nos testes realizados. \\
        \addlinespace
        R02 & Consulta de disponibilidade na unidade
            & O sistema deverá permitir consultar a disponibilidade de um produto na unidade selecionada.
            & A especificação-meta é que pelo menos \textbf{95\% das consultas apresentem informações corretas de disponibilidade}. \\
        \addlinespace
        R03 & Tempo para concluir o pagamento
            & O fluxo de pagamento no aplicativo deverá ser concluído em menos tempo que a espera no caixa.
            & A especificação-meta é que o processo seja concluído em até \textbf{60 segundos}, desconsiderando eventuais atrasos externos relacionados a instituições financeiras ou serviços de terceiros.\\
        \addlinespace
        R04 & Conclusão da compra sem intervenção de funcionário
            & O cliente deverá conseguir identificar a peça, pagar e liberar a etiqueta sem ajuda de um funcionário.
            & A especificação-meta é que pelo menos \textbf{80\% dos usuários concluam a jornada de forma autônoma}. \\
        \addlinespace
        R05 & Clareza do fluxo no aplicativo
            & O fluxo principal do aplicativo deverá poder ser concluído sem instruções adicionais.
            & A especificação-meta é que pelo menos \textbf{85\% dos participantes concluam corretamente o fluxo principal} na primeira tentativa. \\
        \addlinespace
        R06 & Liberação condicionada ao pagamento
            & A etiqueta deverá liberar a peça somente após a confirmação válida do pagamento, e a confirmação exibida ao cliente deverá corresponder ao estado real da transação.
            & A especificação-meta é atingir \textbf{90\% de correspondência entre pagamento confirmado, confirmação apresentada ao usuário e autorização para liberação} nos testes controlados. \\
        \addlinespace
        R07 & Acesso a atendimento humano
            & Todo cenário de falha previsto no protótipo deverá oferecer uma alternativa de atendimento humano que permita concluir a compra.
            & A especificação-meta é que \textbf{90\% dos cenários de falha previstos no protótipo possuam uma alternativa de atendimento humano}, permitindo a continuidade ou conclusão da compra. \\
        \bottomrule
    \end{tabularx}
    \source{Elaborado pelos autores (2026).}
\end{table}

Os requisitos cobrem as dimensões presentes nas vozes selecionadas. O apoio à localização (R01) e a consulta de disponibilidade (R02) respondem às dificuldades na busca pelos produtos e à necessidade de confiar nas informações do aplicativo. A localização foi formulada em nível amplo, e sua forma de implementação ficou para uma etapa posterior, porque as evidências apontam para a dificuldade de encontrar e confirmar a existência da peça, sem indicar que uma localização precisa dentro da loja seja indispensável. O tempo de pagamento (R03) e a conclusão sem intervenção (R04) atendem à demanda por uma jornada mais rápida e menos dependente do caixa, e a clareza do fluxo (R05) está associada à facilidade de uso da tecnologia.

A liberação condicionada ao pagamento (R06) trata ao mesmo tempo da confiança do consumidor e da redução de verificações na saída, uma vez que o sistema deverá fornecer evidência clara de que a transação foi concluída e permitir a saída sem insegurança. Por fim, o acesso a atendimento humano (R07) garante que a automação não elimine o suporte ao consumidor, aspecto relevante sobretudo para os clientes que valorizam assistência durante a jornada. Esse é também o único requisito cuja meta foi fixada em cobertura total sem depender de medição, já que decorre de uma decisão de projeto.

\subsection{Matriz QFD}
\label{subsec:QFD}

Na matriz de relacionamentos do QFD, cada voz do cliente foi relacionada a cada requisito técnico em quatro intensidades: ausência de relação (0), relação fraca (1), moderada (3) e forte (9). A importância técnica de cada requisito resultou da soma dos produtos entre a importância da voz, em escala de 1 a 5, e a intensidade de sua relação com o requisito, o que permite identificar quais características do produto têm maior impacto potencial sobre a satisfação dos clientes. A matriz completa está na Tabela~\ref{tab:qfd}.

\begin{table}[htbp]
    \centering
    \caption{Matriz de relacionamentos do QFD}
    \label{tab:qfd}
    \footnotesize
    \setlength{\tabcolsep}{3.5pt}
    \renewcommand{\arraystretch}{1.25}
    \begin{tabularx}{\textwidth}{@{}l >{\RaggedRight\arraybackslash}X c *{7}{c} c@{}}
        \toprule
        \textbf{Código} & \textbf{Voz do cliente} & \textbf{Imp.} & \textbf{R01} & \textbf{R02} & \textbf{R03} & \textbf{R04} & \textbf{R05} & \textbf{R06} & \textbf{R07} & \textbf{Total} \\
        \midrule
        N03  & Encontrar o tamanho, a cor ou o modelo procurado dentro de um esforço razoável & 3 & 9 & 1 & 0 & 0 & 0 & 0 & 0 & 30 \\
        N05  & Confiar que a informação de estoque digital corresponde à loja física & 2 & 0 & 9 & 0 & 0 & 0 & 0 & 0 & 18 \\
        N06  & Concluir a compra sem risco de ser mal interpretado ou constrangido na saída & 5 & 0 & 0 & 1 & 3 & 3 & 9 & 3 & 95 \\
        N08  & Preservar a opção de atendimento humano mesmo havendo autoatendimento & 1 & 0 & 0 & 1 & 3 & 3 & 1 & 9 & 17 \\
        D03  & Ter confirmação clara e imediata de que o pagamento foi processado & 4 & 0 & 0 & 0 & 0 & 1 & 9 & 3 & 52 \\
        DE01 & Reduzir o tempo de espera no caixa em períodos de alta demanda & 5 & 0 & 0 & 9 & 1 & 1 & 9 & 0 & 100 \\
        DE05 & Ter opção de pagamento rápida para compras simples, sem abrir mão de atendimento nas elaboradas & 4 & 0 & 0 & 9 & 9 & 9 & 3 & 3 & 132 \\
        \midrule
        \multicolumn{3}{@{}l}{Importância técnica} & 27 & 21 & 87 & 59 & 63 & 139 & 48 & 444 \\
        \multicolumn{3}{@{}l}{Importância relativa (\%)} & 6,1 & 4,7 & 19,6 & 13,3 & 14,2 & 31,3 & 10,8 & 100 \\
        \bottomrule
    \end{tabularx}
    \source{Elaborado pelos autores (2026).}
\end{table}

\textcolor{red}{REMOVER ESSA TABELA ACIMA E COLOCAR O QFD DO ZANCUL}


A importância técnica concentrou-se em poucos requisitos. A liberação condicionada ao pagamento (R06) obteve o maior valor, 139 pontos, cerca de 31,3\% do total, principalmente pelas relações fortes com a conclusão da compra sem constrangimento na saída (N06) e com a redução da espera no caixa (DE01), além das relações com a confirmação do pagamento (D03) e com a opção de pagamento rápido (DE05). Garantir que o sistema reconheça corretamente o pagamento e libere a peça constituiu, assim, o principal foco técnico indicado pela matriz.

Em seguida, o tempo para concluir o pagamento (R03) somou 87 pontos, cerca de 19,6\% do total, associado sobretudo à redução da espera no caixa e à opção de pagamento rápido para compras simples. A clareza do fluxo no aplicativo (R05) somou 63 pontos (14,2\%), com relações fortes ou moderadas com a autonomia, a confiança na saída da loja e a compra rápida, o que mostra que disponibilizar o pagamento pelo aplicativo não basta se o usuário não conseguir executar o fluxo corretamente. A conclusão da compra sem intervenção de funcionário (R04) somou 59 pontos (13,3\%), com contribuição concentrada na possibilidade de comprar rapidamente sem depender do atendimento convencional.

O acesso a atendimento humano (R07) somou 48 pontos (10,8\%), vindos principalmente da necessidade de preservar essa opção mesmo havendo autoatendimento (N08). A posição intermediária indica que o atendimento humano continua relevante para a solução, com peso técnico agregado menor que o dos atributos ligados ao pagamento e à conclusão da jornada. O apoio à localização da peça (R01) somou 27 pontos (6,1\%), concentrados na voz sobre encontrar o tamanho, a cor ou o modelo procurado (N03), e a consulta de disponibilidade (R02), 21 pontos (4,7\%), concentrados na confiança na informação de estoque digital (N05).

Os três requisitos de maior importância, R06, R03 e R05, somaram cerca de 65,0\% do total, o que concentrou as prioridades de desenvolvimento na confiabilidade da transação, na rapidez do pagamento e na execução correta do fluxo pelo usuário. Os requisitos ligados à busca e à disponibilidade ficaram com importância menor, embora estejam associados a necessidades relevantes dos clientes, porque a importância técnica depende também do número e da intensidade das relações de cada requisito: as vozes N03 e N05 relacionam-se essencialmente a um requisito cada, enquanto os requisitos de pagamento se relacionam com mais vozes, algumas de importância elevada. Esse resultado reforçou a decisão de manter a localização da peça como requisito secundário, a detalhar depois que o fluxo de pagamento e liberação estiver definido.

A correlação entre os próprios requisitos está na Tabela~\ref{tab:teto_qfd}. Valores positivos indicam que melhorar um requisito tende a favorecer o outro, e valores negativos, que tende a prejudicá-lo.

\begin{table}[htbp]
    \centering
    \caption{Correlação entre os requisitos técnicos (teto da casa da qualidade)}
    \label{tab:teto_qfd}
    \footnotesize
    \renewcommand{\arraystretch}{1.2}
    \begin{tabular}{@{}l*{6}{c}@{}}
        \toprule
            & \textbf{R02} & \textbf{R03} & \textbf{R04} & \textbf{R05} & \textbf{R06} & \textbf{R07} \\
        \midrule
        \textbf{R01} & 3 & 0 & 0 & 0 & 0 & 0 \\
        \textbf{R02} &   & 0 & 0 & 1 & 0 & 0 \\
        \textbf{R03} &   &   & 3 & 3 & 1 & $-1$ \\
        \textbf{R04} &   &   &   & 3 & 1 & $-1$ \\
        \textbf{R05} &   &   &   &   & 3 & $-1$ \\
        \textbf{R06} &   &   &   &   &   & 0 \\
        \bottomrule
    \end{tabular}
    \source{Elaborado pelos autores (2026).}
\end{table}

\textcolor{red}{REMOVER ESSA TABELA ACIMA E COLOCAR O QFD DO ZANCUL}

O tempo de pagamento (R03), a conclusão sem intervenção (R04) e a clareza do fluxo (R05) reforçam-se mutuamente, assim como a localização (R01) e a disponibilidade (R02). O acesso a atendimento humano (R07), por sua vez, apresentou correlação negativa com R03, R04 e R05, o que expressa um conflito de projeto: oferecer atendimento humano durante a compra tende a alongar o fluxo e a reduzir a autonomia. A solução precisará acomodar esse conflito, por exemplo mantendo o atendimento como opção acionada pelo cliente, fora do fluxo principal.

\subsection{Benchmarking competitivo}
\label{subsec:benchmarking}

Para completar a casa da qualidade, a equipe comparou a solução proposta com as alternativas de pagamento e saída já em uso no varejo de moda brasileiro, seguindo a análise do perfil técnico e de mercado recomendada por \citeonline{rozenfeld2006gestao} para a definição das especificações-meta. Foram consideradas quatro alternativas: a rede atual da Riachuelo, em que a maior parte das cerca de 440 lojas ainda opera com caixa tradicional e apenas as lojas conceito contam com autoatendimento; a Renner, com terminais de autoatendimento por radiofrequência em 213 lojas; a C\&A, que passou a oferecer autoatendimento com etiquetas de radiofrequência no conceito de loja Energia, lançado em 2025 \cite{brazileconomy2025cea}; e a Zara, cujas lojas brasileiras reabertas a partir de 2023 incluem leitura digital das peças, provadores automáticos e consulta de estoque \cite{infomoney2023zara}.

A comparação foi feita em duas partes. Na avaliação competitiva, a equipe atribuiu a cada alternativa uma nota de 1 a 5 para o grau em que atende cada voz do cliente selecionada, com base nos recursos documentados em fontes públicas e nas Seções~\ref{sec:mercado} e \ref{sec:analise_clientes}. A escala adotada foi a seguinte: 5 indica recurso documentado que atende plenamente a voz; 4, recurso documentado com alguma limitação ou restrito a parte das lojas; 3, atendimento parcial ou mecanismo não descrito nas fontes; 2, ausência de recurso documentado, com o processo dependente do vendedor ou do caixa; e 1, recurso inexistente. A coluna da solução proposta registra a nota pretendida pela equipe e deverá ser confirmada nos testes do protótipo. As notas estão reunidas na Tabela~\ref{tab:benchmark_competitivo}, em que a média ponderada usa a importância de cada voz na matriz do QFD.

\begin{table}[htbp]
    \centering
    \caption{Avaliação competitiva das alternativas de pagamento e saída da loja}
    \label{tab:benchmark_competitivo}
    \footnotesize
    \setlength{\tabcolsep}{4pt}
    \renewcommand{\arraystretch}{1.25}
    \begin{tabularx}{\textwidth}{@{}l >{\RaggedRight\arraybackslash}X c c c c c c@{}}
        \toprule
        \textbf{Código} & \textbf{Voz do cliente} & \textbf{Imp.} & \textbf{Riachuelo} & \textbf{Renner} & \textbf{C\&A} & \textbf{Zara} & \textbf{Proposta} \\
        \midrule
        N03 & Encontrar o produto procurado com facilidade & 3 & 2 & 2 & 4 & 4 & 3 \\
        N05 & Confiar na informação de estoque digital & 2 & 2 & 2 & 3 & 4 & 4 \\
        N06 & Sair da loja sem constrangimento por falha ou checagem & 5 & 3 & 4 & 3 & 3 & 5 \\
        N08 & Contar com atendimento humano quando necessário & 1 & 5 & 4 & 4 & 3 & 4 \\
        D03 & Receber confirmação imediata do pagamento & 4 & 4 & 4 & 4 & 4 & 5 \\
        DE01 & Reduzir o tempo de espera no caixa & 5 & 2 & 4 & 3 & 4 & 5 \\
        DE05 & Ter pagamento rápido sem perder o atendimento & 4 & 2 & 5 & 4 & 3 & 4 \\
        \midrule
        \multicolumn{3}{@{}l}{Média ponderada} & 2,67 & 3,75 & 3,50 & 3,58 & 4,46 \\
        \bottomrule
    \end{tabularx}
    \source{Elaborado pelos autores (2026) com base em \citeonline{consumidormoderno2025renner}, \citeonline{brazileconomy2025cea}, \citeonline{infomoney2023zara}, \citeonline{ecommercenews2020zara} e \citeonline{tribunapr2026riachuelo}.}
\end{table}

A rede atual da Riachuelo obteve a menor média ponderada (2,67), abaixo da C\&A (3,50), da Zara (3,58) e da Renner (3,75). A diferença concentrou-se nas duas vozes de maior peso ligadas ao pagamento: na redução da espera no caixa (DE01) e na oferta de uma via rápida sem perder o atendimento (DE05), a Riachuelo recebeu nota 2, porque o autoatendimento existe apenas nas lojas conceito, enquanto a Renner já o oferece em mais da metade da rede e mantém, em paralelo, o caixa, o pagamento pelo aplicativo e a venda móvel com funcionário \cite{consumidormoderno2025renner}. Nas vozes de busca e estoque (N03 e N05), a C\&A e a Zara ficaram à frente por oferecerem consulta de tamanhos ou de estoque pelo aplicativo e, no caso da Zara, localização da peça no plano da loja, recurso documentado nas lojas espanholas \cite{ecommercenews2020zara}. No atendimento humano (N08), a Riachuelo recebeu a nota mais alta, já que o caixa tradicional mantém um funcionário em todas as compras.

A solução proposta pretende alcançar média 4,46, com as maiores notas justamente nas vozes em que a Riachuelo ficou mais atrás, a espera no caixa e a saída sem constrangimento, esta última apoiada no indicador luminoso que confirma a liberação da peça. Nas vozes de busca e estoque, a nota pretendida é menor, porque a localização da peça foi tratada como requisito secundário (Seção~\ref{subsec:requisitos_do_produto}).

Na segunda parte, o benchmarking técnico, a equipe registrou como cada alternativa atende cada requisito do produto. Como as fontes públicas descrevem os recursos, mas não publicam tempos de pagamento nem taxas de erro, a característica documentada de cada alternativa está na Tabela~\ref{tab:benchmark_tecnico}, e os valores numéricos serão obtidos nas visitas descritas a seguir.

\begin{table}[htbp]
    \centering
    \caption{Benchmarking técnico por requisito do produto}
    \label{tab:benchmark_tecnico}
    \scriptsize
    \setlength{\tabcolsep}{3pt}
    \renewcommand{\arraystretch}{1.3}
    \begin{tabularx}{\textwidth}{@{}
        >{\RaggedRight\arraybackslash}p{1.9cm}
        >{\RaggedRight\arraybackslash}X
        >{\RaggedRight\arraybackslash}X
        >{\RaggedRight\arraybackslash}X
        >{\RaggedRight\arraybackslash}X
        >{\RaggedRight\arraybackslash}p{1.5cm}@{}}
        \toprule
        \textbf{Requisito} & \textbf{Riachuelo} & \textbf{Renner} & \textbf{C\&A} & \textbf{Zara} & \textbf{Meta} \\
        \midrule
        R01 Localização da peça & N.D. & N.D. & Aplicativo localiza peças (Energia) & Aplicativo mostra a peça no plano da loja (Espanha) & 90\% \\
        \addlinespace
        R02 Disponibilidade & N.D. & N.D. & Tamanhos em tempo real no aplicativo & Consulta de estoque na loja e nos provadores & 95\% \\
        \addlinespace
        R03 Tempo de pagamento & A medir & A medir & A medir; biometria facial para o cartão C\&A Pay & A medir & 60 s \\
        \addlinespace
        R04 Compra sem funcionário & Não, exceto nas lojas conceito & Sim, 742 terminais em 213 lojas & Sim, nas lojas Energia & Leitura digital das peças & 80\% \\
        \addlinespace
        R05 Clareza do fluxo & A medir & A medir & A medir & A medir & 85\% \\
        \addlinespace
        R06 Liberação da peça & Operador remove a etiqueta & Leitura por radiofrequência no pagamento & Radiofrequência nas lojas Energia; retirada pelo cliente nas demais & N.D. & 90\% \\
        \addlinespace
        R07 Atendimento humano & Caixa em todas as compras & Venda móvel com funcionário & Vendedores como curadores; totens nos provadores & N.D. & 90\% \\
        \bottomrule
    \end{tabularx}
    \source{Elaborado pelos autores (2026). N.D.: não documentado em fonte pública.}
\end{table}

As medições serão feitas em visitas às lojas de Riachuelo, Renner, C\&A e Zara, no mesmo dia e horário e, sempre que possível, no mesmo shopping. Em cada loja, duas pessoas da equipe realizarão a mesma compra padrão, de uma peça básica, e registrarão em ficha o tempo de fila, o tempo de pagamento, o número de passos e de erros no processo, o tempo para conseguir ajuda de um funcionário e a forma de liberação da peça. Com esses valores, a equipe deverá revisar as especificações-meta da Tabela~\ref{tab:requisitos}, já que uma meta só representará vantagem competitiva se superar o desempenho medido nas alternativas existentes.

Como a solução é formada pela integração entre a etiqueta e o aplicativo, a equipe concluiu, a partir da matriz, que o desenvolvimento deverá priorizar o funcionamento do sistema integrado, capaz de proporcionar uma jornada de compra rápida, confiável e executável pelo próprio consumidor, mais do que o desempenho isolado de cada componente.
